% ===================================================================
%  ТАБЛИЦЯ № 9 — Гори і курорт. — Ліс і полювання.
%  Traduction 2026 d'apres texto/io, controlee sur texto/fr, texto/en
%  et texto/es. Consigne : voir texto/uk/10-tablycia-01.tex.
%
%  LE LEXIQUE DE LA MONTAGNE ET DU BOIS. La mine est une « копальня »
%  et non un « рудник », la fonderie une « ливарня » et non un
%  « литейный завод », la carriere une « каменярня », le bucheron un
%  « дроворуб », le charbonnier un « вугляр » : l'ukrainien nomme le
%  metier par le verbe de la chose faite, la ou le russe emprunte plus
%  volontiers au vocabulaire administratif.
% ===================================================================

%%K t09-apar-1 apar
\VUsaut{4.0mm}
\VUtitre{15.0pt}{60}{ТАБЛИЦЯ № 9}
\VUsaut{3.0mm}

%%K t09-tit-1 sub
\VUcentre{12.0pt}{0}{\textit{Перша сцена.}}
\VUsaut{3.0mm}

%%K t09-tit-2 sub
\VUpk{11.4pt}{40}{Гори. --- Курортне містечко.}
\VUsaut{3.0mm}

%%K t09-c1-01-1 p
1. --- Я вже тиждень у Праці. Не можу переказати всього подиву, якого я зазнав,
ставши перед чудесним \VUgras{гірським хребтом} \textsuperscript{(1)} Піренеїв,
чий \VUgras{гребінь} \textsuperscript{(2)} вирізується на синьому небі, як
велетенська пилка.

%%K t09-c1-02-1 p
2. --- Я не уявляв, що піренейські \VUgras{вершини} \textsuperscript{(3)} можуть
сягати такої \VUgras{висоти}, і зостався вражений чудовими \VUgras{льодовиками}
\textsuperscript{(5)} і сніговими \VUgras{шпилями} \textsuperscript{(4)}, по яких
мчать такі страшні \VUgras{лавини}.

%%K t09-tit-3 sub
\VUsaut{3.0mm}
\VUpk{10.2pt}{40}{Гірський краєвид.}
\VUsaut{3.0mm}

%%K t09-c1-03-1 p
3. --- Наступного ж дня по приїзді я пішов до старого згаслого \VUgras{вулкана}
\textsuperscript{(6)} біля Праца. На голому безплідному краю \VUgras{кратера} ще
ясно видно сліди, лишені \VUgras{лавою}, яка текла кілька віків тому по
\VUgras{схилу гори} \textsuperscript{(7)}. На одному з \VUgras{укосів} мені
пощастило знайти кілька шматків цієї лави.

%%K t09-c1-04-1 p
4. --- Прац стоїть на кордоні, і \VUgras{фортеця} \textsuperscript{(8)}, що
стереже \VUgras{перевал} \textsuperscript{(27)} між Францією та Іспанією, дуже
схожа на ті старі замки понад берегами Рейну, які немов виглядають
\VUgras{здобич} із вершини своєї скелі.

%%K t09-c1-05-1 p
5. --- Величезний \VUgras{орел} \textsuperscript{(9)}, у якого гніздо недалеко
від \VUgras{форту} \textsuperscript{(8)}, часто привертає мою увагу. Цей
\VUgras{хижий птах} з могутнім \VUgras{дзьобом} часто приносить пташенятам
ягня або козеня, затиснене в \VUgras{пазурах}. \VUgras{Крила} в нього такі
великі, що він довго може ширяти зі своєю важкою ношею.

%%K t09-tit-4 sub
\VUsaut{3.0mm}
\VUpk{10.2pt}{40}{Пішохід.}
\VUsaut{3.0mm}

%%K t09-c1-06-1 p
6. --- Найприємніша тутешня прогулянка --- похід до \VUgras{водоспаду}
\textsuperscript{(10)} Ко. \VUgras{Падуча вода} валиться вихором і потім зникає
гарненьким \VUgras{струмком} у \VUgras{ялиновому лісі} \textsuperscript{(11)} на
захід від містечка. Ми піднялися цим лісом до \VUgras{метеорологічної станції}
\textsuperscript{(12)} і з вершини \VUgras{вежі} охопили поглядом усю
\VUgras{панораму} округи. \VUgras{Краєвид} чудовий, і \VUgras{природа},
здається, розточила цьому краєві всі скарби, які були в її розпорядженні.

%%K t09-c1-07-1 p
7. --- Щоб спуститися, ми взяли \VUgras{фунікулер} \textsuperscript{(13)} і
зійшли коло маленької \VUgras{станції} біля \VUgras{руїн} \textsuperscript{(14)}
старого \VUgras{феодального замку}, у якому колись жили володарі Праца.

%%K t09-c1-08-1 p
8. --- Бідний замок! Що має він думати, бачачи внизу під собою охайне
\VUgras{курортне містечко} \textsuperscript{(15)}, яке стало модним
\VUgras{водним курортом}?

%%K t09-c1-08-2 p
\VUgras{Клімат} тут такий приємний, а \VUgras{мінеральна вода} така цілюща, що
\VUgras{курс лікування}, який проходять у \VUgras{санаторії}
\textsuperscript{(17)}, --- справжня втіха.

%%K t09-tit-5 sub
\VUsaut{3.0mm}
\VUpk{10.2pt}{40}{Приємне курортне містечко.}
\VUsaut{3.0mm}

%%K t09-c1-09-1 p
9. --- І справді, зроблено все, щоб перебування було приємне. Збудували великий
\VUgras{віадук} \textsuperscript{(16)} для залізниці; \VUgras{парк}
\textsuperscript{(18)} при \VUgras{курзалі}, де дають \VUgras{концерти},
розпланований чарівно. Кілька зграбних \VUgras{вілл} \textsuperscript{(19)}
збудували навколо \VUgras{казино} \textsuperscript{(21)}, а дуже добре доглянуте
\VUgras{шосе} \textsuperscript{(22)} робить пересування надзвичайно легким.

%%K t09-c1-10-1 p
10. --- Простий \VUgras{насип} \textsuperscript{(23)} відділяє \VUgras{дорогу}
\textsuperscript{(20)} від самої \VUgras{долини} \textsuperscript{(25)}, на дні
якої лежить граційне \VUgras{озерце} \textsuperscript{(26)}, що живить прозорий
\VUgras{струмок} \textsuperscript{(24)}.

%%K t09-c1-11-1 p
11. --- По другий бік долини розробляють залізну \VUgras{копальню}
\textsuperscript{(28)}. Я кілька разів ходив дивитися, як \VUgras{рудокопи}
спускаються в \VUgras{шахту}, і навіть входив у \VUgras{штольню}, де вони
проводять день, добуваючи \VUgras{руду}, у якій міститься \VUgras{метал}.

%%K t09-c1-12-1 p
12. --- Біля копальні збудували велику \VUgras{ливарню}, щоб тут же вжити
видобуті з неї багатства. \VUgras{Домна} \textsuperscript{(29)}, топлена
тутешнім рясним \VUgras{вугіллям}, дозволяє добувати \VUgras{чавун} швидко і
дешево.

%%K t09-c1-13-1 p
13. --- Недалеко від копальні розробляють \VUgras{каменярню}
\textsuperscript{(30)}. Старий \VUgras{каменяр} відколює своєю \VUgras{киркою}
не \VUgras{граніт}, не \VUgras{порфір} і навіть не \VUgras{пісковик}, а простий
\VUgras{будівельний камінь} на будову домів.

%%K t09-c1-14-1 p
14. --- Одна з найкращих оздоб цього краю --- \VUgras{ялина}
\textsuperscript{(31)}. Скрізь --- на дні \VUgras{яру} \textsuperscript{(32)}, на
березі \VUgras{потоку} \textsuperscript{(33)}, коло \VUgras{провалля}
\textsuperscript{(34)} чи урвища --- її тонкі голки вносять свою зелену ноту.

%%K t09-tit-6 sub
\VUsaut{3.0mm}
\VUpk{10.2pt}{40}{Прогулянки пішки.}
\VUsaut{3.0mm}

%%K t09-c1-15-1 p
15. --- Я користаюся з канікул, щоб чинити довгі прогулянки. \VUgras{Стежки}
\textsuperscript{(35)}, що в'ються по горі, дозволяють пройти кілька
\VUgras{кілометрів}, а часто й кілька \VUgras{верст}, не помічаючи відстані.

%%K t09-c1-15-2 p
\VUgras{Верстові стовпи} \textsuperscript{(36)} і \VUgras{покажчики}
\textsuperscript{(37)} розмічають стежки, на яких часто стрічаєш молодого
\VUgras{горянина} \textsuperscript{(38)} з \VUgras{торбиною}
\textsuperscript{(40)} на боці, що несе \VUgras{бабака} \textsuperscript{(39)},
або якогось \VUgras{цигана} \textsuperscript{(41)}, у якого \VUgras{хутряна
шапка} \textsuperscript{(42)} дивно сперечається зі старим подертим
\VUgras{плащем} \textsuperscript{(43)}. Він веде на \VUgras{ланцюгу} вченого
\VUgras{ведмедя} \textsuperscript{(44)}.

%%K t09-tit-7 sub
\VUpk{10.2pt}{40}{Сходження.}
\VUsaut{3.0mm}

%%K t09-c1-16-1 p
16. --- Майже щодня я ходжу з \VUgras{провідником} \textsuperscript{(48)} на
\VUgras{сходження}, і я дуже гордий, коли з \VUgras{ранцем}
\textsuperscript{(47)} за плечима і \VUgras{альпенштоком} у руці, як справжній
\VUgras{турист} \textsuperscript{(46)}, беру приступом \VUgras{скелі}
\textsuperscript{(45)}.

%%K t09-c1-17-1 p
17. --- З вершини гори люди на рівнині здаються не більші за мурах;
\VUgras{отара овець} \textsuperscript{(50)}, що пасеться на \VUgras{пасовиську}
\textsuperscript{(49)}, схожа на дитячу забавку. \VUgras{Сосна}
\textsuperscript{(51)} чи навіть дерев'яне \VUgras{шале} \textsuperscript{(52)}
--- ледве помітна цятка на зелені.

%%K t09-c1-18-1 p
18. --- На цих прогулянках ми часто стрічаємо на \VUgras{стежці}
\textsuperscript{(53)} \VUgras{мисливця на серн} \textsuperscript{(54)}, що несе
на плечі щойно вбитого звіра.

%%K t09-c1-19-1 p
19. --- Я поклав у свій гербарій дуже примітний лист \VUgras{папороті}
\textsuperscript{(55)} і гарненьку рожеву гілочку \VUgras{вересу}
\textsuperscript{(57)}, яку зірвав коло входу до маленької \VUgras{печери}
\textsuperscript{(56)} на останній прогулянці.

%%K t09-tit-8 sub
\VUsaut{3.0mm}
\VUcentre{12.0pt}{0}{\textit{Друга сцена.}}
\VUsaut{3.0mm}

%%K t09-tit-9 sub
\VUpk{11.4pt}{40}{Ліс. --- Полювання.}
\VUsaut{3.0mm}

%%K t09-c2-01-1 p
1. --- Учора я провів у лісі чудовий день. Діяльне життя звірів і \VUgras{комах}
\textsuperscript{(5)}, які там живуть, дуже мене займало.

%%K t09-c2-01-2 p
\VUgras{Павук} \textsuperscript{(1)}, що тче свою \VUgras{павутину}
\textsuperscript{(2)} між двома гілками, щоб піймати пролітну \VUgras{муху}
\textsuperscript{(3)}; \VUgras{слимак} \textsuperscript{(4)}, що важко тягнеться
зі своїм будиночком; жук-олень, що виглядає здобич; клопітка мураха, що працює
коло \VUgras{мурашника} \textsuperscript{(6)} --- усі ці малі створіння снували й
трудилися з незвичайною силою.

%%K t09-c2-02-1 p
2. --- \VUgras{Змія} \textsuperscript{(7)}, яку я спершу взяв за \VUgras{гадюку},
а вона виявилася безневинним \VUgras{вужем}, згорнулася під стовбуром дерева і
час від часу вистромляла свою тонку голівку, завжди немов готову вжалити. Передо
мною важко повз великий \VUgras{слизняк} \textsuperscript{(8)}, який з'їв одну зі
смачних маленьких \VUgras{суниць} \textsuperscript{(11)}, яких тут безліч.

%%K t09-c2-03-1 p
3. --- Земля була встелена \VUgras{лісовими квітами}: \VUgras{первоцвіти}
\textsuperscript{(9)}, \VUgras{барвінки} \textsuperscript{(10)} і багато інших
невідомих мені рослин цвіли серед \VUgras{кущиків трави}
\textsuperscript{(12)}.

%%K t09-c2-04-1 p
4. --- У цьому краю \VUgras{трюфель} майже такий самий звичайний, як
\VUgras{гриб} \textsuperscript{(14)}, і \VUgras{порося} \textsuperscript{(13)}, що
належить одному лісникові, жадібно рилося, чи не знайдеться їх.

%%K t09-tit-10 sub
\VUsaut{3.0mm}
\VUpk{10.2pt}{40}{Браконьєрство.}
\VUsaut{3.0mm}

%%K t09-c2-05-1 p
5. --- Я застав \VUgras{кінець} сцени, яка мене дуже потішила. Завдяки чуттю
своєї \VUgras{такси} \textsuperscript{(15)} \VUgras{лісник} \textsuperscript{(16)}
щойно накрив \VUgras{браконьєра} \textsuperscript{(22)} тієї хвилини, коли той
збирався винести вбиту ним \VUgras{дичину} \textsuperscript{(17)}. Волівши
розлучитися зі здобиччю, аби не бути затриманим, браконьєр утік, і
\VUgras{заєць} \textsuperscript{(18)}, \VUgras{лань} \textsuperscript{(19)},
\VUgras{козуля} \textsuperscript{(20)} і \VUgras{кабан} \textsuperscript{(21)}
зосталися лежати на траві, на превелику радість лісникові.

%%K t09-c2-06-1 p
6. --- Вражає розмаїття дерев у цьому лісі. \VUgras{Буки} \textsuperscript{(23)}
і \VUgras{берези} \textsuperscript{(26)}, \VUgras{сосни}, що дають \VUgras{смолу}
\textsuperscript{(27)}, \VUgras{дуби} \textsuperscript{(29)} і багато інших
переплітають своє гілля.

%%K t09-c2-06-2 p
\VUgras{Плющ} \textsuperscript{(24)}, що чіпляється за стовбури високих дерев,
надає \VUgras{високостовбурному лісові} \textsuperscript{(25)} яскравої зелені,
яка приємно зливається з блідим і м'яким відтінком \VUgras{моху}
\textsuperscript{(31)}, де \VUgras{зелений дятел} \textsuperscript{(30)} шукає
комах.

%%K t09-tit-11 sub
\VUsaut{3.0mm}
\VUpk{10.2pt}{40}{Лісовий краєвид.}
\VUsaut{3.0mm}

%%K t09-c2-07-1 p
7. --- Старі дуби мене зачарували. Їхнє велике вузлувате \VUgras{коріння}
\textsuperscript{(32)}, наполовину вийшле із землі, надає їм чудового вигляду
сили і моці, і я не розумію, як \VUgras{власник} \textsuperscript{(35)} може
зважитися їх валити і віддавати \VUgras{пилці} \textsuperscript{(34)}
\VUgras{дроворуба} \textsuperscript{(33)}. Бідні \VUgras{стовбури}
\textsuperscript{(36)}, позбавлені \VUgras{кори} \textsuperscript{(37)} або
розколоті \VUgras{сокирою} \textsuperscript{(38)} \VUgras{поденника}
\textsuperscript{(39)}, засмучують мене.

%%K t09-c2-08-1 p
8. --- Поруч старий \VUgras{вугляр} \textsuperscript{(41)} склав своєю
\VUgras{лопатою} \VUgras{купу} \textsuperscript{(42)} вугілля, а стара жінка
несла \VUgras{в'язку} \textsuperscript{(40)} \VUgras{хмизу} з гілок зрубаних
дерев.

%%K t09-tit-12 sub
\VUsaut{3.0mm}
\VUpk{10.2pt}{40}{Лови.}
\VUsaut{3.0mm}

%%K t09-c2-09-1 p
9. --- Я збирався трохи спочити в \VUgras{домі лісника} \textsuperscript{(46)},
але, дійшовши до маленького \VUgras{ставу} \textsuperscript{(43)}, по якому
плавав гарний \VUgras{лебідь} \textsuperscript{(45)}, я помітив, що недавня
\VUgras{повінь} \textsuperscript{(44)} обернула стежку на справжнє
\VUgras{болото}. Довелося зробити гак, і, проходячи повз \VUgras{кедр}
\textsuperscript{(49)}, \VUgras{тополі} \textsuperscript{(48)} і \VUgras{в'яз}
\textsuperscript{(47)}, що стоять на узліссі, я побачив, як із \VUgras{хащів}
\textsuperscript{(54)} вирвався чудовий \VUgras{олень} \textsuperscript{(50)},
переслідуваний невтомною \VUgras{зграєю} \textsuperscript{(51)}, яку гарячив
\VUgras{ріг} \textsuperscript{(53)} \VUgras{кінних мисливців}
\textsuperscript{(52)}. Бідний звір зовсім вибився з сил. Він упав, конаючи, за
кілька кроків від мене.

%%K t09-c2-10-1 p
10. --- Син лісника, приваблений гамором \VUgras{ловів}, вийшов з дому, і ми
вдвох вернулися в ліс. Він примітив \VUgras{сороку} \textsuperscript{(56)} на
гілці старого дуба. Він думав, що гніздо її сховане в \VUgras{листі}
\textsuperscript{(58)}, і хотів його взяти.

%%K t09-c2-10-2 p
Меткий, як \VUgras{білка} \textsuperscript{(61)}, він ліз із \VUgras{гілки}
\textsuperscript{(55)} на гілку, але нічого не знайшов і тільки сполохав стару
\VUgras{ворону} \textsuperscript{(60)}, що сиділа на \VUgras{суку}
\textsuperscript{(57)}.
\VUsaut{4.0mm}
\VUfilet{22mm}
